% !TEX program = xelatex
% ============================================================
% SLB 技术文档：6.5.7–6.5.9 比特加扰、SC-FDMA、ZC 序列
% 规范：slb-doc-generator / tech-doc-style-guide.md
% 模块类型：通用功能
% ============================================================
\documentclass[11pt,a4paper]{ctexart}

\usepackage{amsmath,amssymb,amsfonts}
\usepackage{booktabs}
\usepackage{geometry}
\usepackage{hyperref}
\usepackage{enumitem}
\usepackage{xcolor}
\usepackage{longtable}
\usepackage{array}

\geometry{left=2.5cm,right=2.5cm,top=2.5cm,bottom=2.5cm}
\hypersetup{colorlinks=true,linkcolor=blue,urlcolor=blue}

\newcolumntype{L}[1]{>{\raggedright\arraybackslash}p{#1}}

\title{%
  \textbf{星闪 SLB 物理层}\\[0.4em]
  \Large 6.5.7\textasciitilde 6.5.9 比特加扰、波形生成与序列技术文档\\[0.3em]
  \large 比特加扰 \textbullet\ SC-FDMA \textbullet\ Zadoff-Chu \textbullet\ 组跳
}
\author{依据 T/XS 10001-2025《星闪无线通信系统 接入层\\
同步低时延宽带空口 SLB 技术要求和测试方法》整理\\
（团体标准 20250326 版）}
\date{2026 年 7 月 27 日}

\begin{document}
\maketitle
\tableofcontents
\newpage

%======================================================================
\part{总览}
%======================================================================

\section{文档范围与模块划分}

本文档对应 T/XS 10001-2025 第 6.5.7\textasciitilde 6.5.9 节，覆盖比特加扰、SC-FDMA（DFT 预编码）与 Zadoff-Chu 参考序列（含组跳）。三节为可复用算法块；Gold 引擎统一复用 6.5.2。

\begin{longtable}{L{2.4cm}L{3.8cm}L{7.5cm}}
\caption{文档范围与模块划分\label{tab:6579-scope}} \\
\toprule
\textbf{子节号} & \textbf{主题} & \textbf{核心输出/行为} \\
\midrule
6.5.7 & 比特加扰 & $\tilde{t}_i=(t_i+c(i))\bmod 2$；$c_{\mathrm{init}}$ 按业务选取 \\
6.5.8 & SC-FDMA 波形 & $M$ 点 DFT 预编码，$1/\sqrt{M}$ 归一化；$L=M_{\mathrm{symb}}/M$ \\
6.5.9 & ZC 参考序列 & $r_{u,v}^{(\alpha)}(n)$；$m{=}2$ 查表，其余按标准闭式 \\
6.5.9.1 & 序列组跳 & $u=(f_{\mathrm{gh}}(n_S)+f_{\mathrm{SS}})\bmod 30$ \\
\bottomrule
\end{longtable}

\section{与相关章节的关系}

\begin{longtable}{L{2.2cm}L{1.8cm}L{9.5cm}}
\caption{与相关章节的关系\label{tab:6579-rel}} \\
\toprule
\textbf{相关节号} & \textbf{关系} & \textbf{说明} \\
\midrule
6.5.2 & 上游 & Gold 生成 $c(i)$；加扰与组跳使用\textbf{不同} $c_{\mathrm{init}}$ \\
6.2.5/6.2.6 & 上游 & 编码/速率匹配比特进 6.5.7 \\
6.5.6 & 上/下游 & 加扰后调制；调制符号进 6.5.8 \\
6.3.3 & 下游 & 深覆盖：加扰$\rightarrow$调制$\rightarrow$SC-FDMA；DMRS 用 6.5.9 \\
6.5.3 & 平行 & RS QPSK 的 $c_{\mathrm{init}}$ 与加扰/组跳三者互不混用 \\
6.2.3 & 下游 & RE 映射 \\
\bottomrule
\end{longtable}

\section{通用功能端到端流水线}

\begin{center}
\small
\textbf{6.2.5 数据}：编码 $\rightarrow$ \textbf{6.5.7} $\rightarrow$ 6.5.6 $\rightarrow$ 映射（无 6.5.8）\\[0.3em]
\textbf{6.3.3 深覆盖}：CRC+\textbf{6.5.7} $\rightarrow$ 6.5.6 $\rightarrow$ \textbf{6.5.8} $\rightarrow$ 映射；\\
并行 DMRS：\textbf{6.5.9.1} 组跳 $\rightarrow$ \textbf{6.5.9} $\rightarrow$ OCC$\times\beta_{\mathrm{RS}}$ 映射
\end{center}

\newpage
%======================================================================
\part{6.5.7\textasciitilde 6.5.9 工程解读}
%======================================================================

\section{协议章节对应}

本节对应 T/XS 10001-2025 第 6 章第 6.5.7、6.5.8、6.5.9（含 6.5.9.1）节。

\section{功能定义}

\begin{enumerate}[nosep]
  \item \textbf{加扰}：用 Gold $c(i)$ 对编码比特逐位模 2 加，随机化比特统计；
  \item \textbf{SC-FDMA 预编码}：对调制符号分段做 $M$ 点 DFT，降低 PAPR；
  \item \textbf{生成 ZC/表序列}：按 $u,v,\alpha,m$ 生成 DMRS 基序列并施加循环移位；
  \item \textbf{组跳}：按时隙更新组索引 $u\in\{0,\ldots,29\}$。
\end{enumerate}

\section{模块边界（上下游及职责划分）}

\subsection{上游模块}

\begin{longtable}{L{4.2cm}L{4.5cm}L{5.5cm}}
\caption{上游模块输入\label{tab:6579-up}} \\
\toprule
\textbf{上游} & \textbf{输入} & \textbf{说明} \\
\midrule
编码/速率匹配 & $t_i$/$g_k$ & 进 6.5.7 \\
调度/RRC & $N_{\mathrm{G-PID}}$、帧号、$n_{\mathrm{ID}}^X$ & $c_{\mathrm{init}}$/组跳 \\
6.5.2 Gold & $c(i)$ & 加扰与组跳 \\
6.5.6 调制 & $x(n)$ & 进 6.5.8 \\
DMRS 调度 & $m,\alpha,n_S$ & 进 6.5.9 \\
\bottomrule
\end{longtable}

\subsection{下游模块}

\begin{longtable}{L{4.2cm}L{4.5cm}L{5.5cm}}
\caption{下游模块输出\label{tab:6579-down}} \\
\toprule
\textbf{下游} & \textbf{输出} & \textbf{说明} \\
\midrule
6.5.6 调制 & $\tilde{t}_i$ & 加扰后比特 \\
资源映射 & $y(n)$、$r_{u,v}^{(\alpha)}$ & SC-FDMA / DMRS \\
接收解扰 & 恢复 $t_i$ & 同 $c_{\mathrm{init}}$ 再异或 \\
\bottomrule
\end{longtable}

\subsection{本模块不负责的内容}

\begin{longtable}{L{5.5cm}L{8.5cm}}
\caption{本模块不负责的内容\label{tab:6579-out}} \\
\toprule
\textbf{不负责} & \textbf{负责节} \\
\midrule
Gold LFSR 内部递推细节 & 6.5.2（本模块只传 $c_{\mathrm{init}}$） \\
星座映射 & 6.5.6 \\
IFFT/CP/上变频 & 6.3.3.6 \\
OCC/$β_{\mathrm{RS}}$ 最终 RE 填充细则 & 6.3.3.2.2 \\
\bottomrule
\end{longtable}

\section{在 SLB 通信流程中的角色}

6.5.7 位于编码之后、调制之前；6.5.8 仅深覆盖路径；6.5.9 与业务数据并行生成 DMRS。接收解扰在译码前，使用发送侧同一 $c_{\mathrm{init}}$。

%----------------------------------------------------------------------
\section{关键参数与强制要求}
%----------------------------------------------------------------------

\subsection{6.5.7 比特加扰}

\begin{equation}
\tilde{t}_i=(t_i+c(i))\bmod 2
\end{equation}

\begin{longtable}{L{4.5cm}L{5.5cm}L{4cm}}
\caption{加扰 $c_{\mathrm{init}}$ 选用\label{tab:6579-cinit}} \\
\toprule
\textbf{业务} & \textbf{$c_{\mathrm{init}}$} & \textbf{标准} \\
\midrule
第一/二类数据（6.2.5） & $n_f\cdot 2^{24}+N_{\mathrm{G-PID}}$ & 6.2.5 \\
系统广播 & $N_{\mathrm{ID}}^{\mathrm{TS}}$ & 6.3.3.3 \\
控制信息（G/T） & $510$（固定） & 6.3.3.4 \\
深覆盖 PHY 数据 & 高层配置 & 6.3.3.5 \\
\bottomrule
\end{longtable}

其中数据默认式中 $n_f$ 为起始无线帧号低 7 位（$\bmod 128$）。

\textbf{三场景 Gold 隔离}（禁止混流）：

\begin{longtable}{L{4cm}L{2.2cm}L{7.5cm}}
\caption{三类 Gold $c_{\mathrm{init}}$ 场景\label{tab:6579-gold3}} \\
\toprule
\textbf{场景} & \textbf{节} & \textbf{$c_{\mathrm{init}}$} \\
\midrule
伪随机 QPSK RS & 6.5.3 & $n,l,N_{\mathrm{ID}},N_{\mathrm{CP}}$ 组合 \\
比特加扰 & 6.5.7 & 上表 \\
ZC 组跳 & 6.5.9.1 & $\lfloor n_{\mathrm{ID}}^X/30\rfloor$ \\
\bottomrule
\end{longtable}

\textbf{金标}：$N_{\mathrm{G-PID}}{=}42$，帧号 $127$ $\Rightarrow$ $n_f{=}127$，$c_{\mathrm{init}}{=}127\cdot 2^{24}+42{=}2130700346$。

\subsection{6.5.8 SC-FDMA}

\begin{align}
M_{\mathrm{SC}}^{\mathrm{Total}} &= M_{\mathrm{SCG},1}\cdot N_{\mathrm{SCG},1}\\
M_{\mathrm{SCG},1} &= 2^{\alpha_2}3^{\alpha_3}5^{\alpha_5}\\
L &= M_{\mathrm{symb}}/M_{\mathrm{SC}}^{\mathrm{Total}}
\end{align}
$M_{\mathrm{symb}}$ 必须被 $M_{\mathrm{SC}}^{\mathrm{Total}}$ 整除。对段 $l$：
\begin{equation}
y(lM+k)=\frac{1}{\sqrt{M}}\sum_{i=0}^{M-1}x(lM+i)\,e^{-j2\pi ik/M}
\end{equation}
$1/\sqrt{M}$ 强制；段间不交叉混合。

\subsection{6.5.9 Zadoff-Chu 与组跳}

\begin{equation}
r_{u,v}^{(\alpha)}(n)=e^{j\alpha n}\,r_{u,v}(n),\quad 0\le n<M_{\mathrm{SC}}^{\mathrm{RS}}
\end{equation}
$u\in\{0,\ldots,29\}$；$1\le m\le 5$ 时仅 $v{=}0$；更大 $m$ 时 $v\in\{0,1\}$。

\textbf{分支}：
\begin{itemize}[nosep]
  \item $m{=}2$：标准表（相位 $\varphi(n)\in\{1,-1,3,-3\}$），$r_{u,0}(n)=\exp(-j\varphi(n)\pi/4)$；\textbf{禁止}用 ZC 闭式；
  \item 其他长度：按标准取最大素数 $N_{\mathrm{ZC}}^{\mathrm{RS}}<M_{\mathrm{SC}}^{\mathrm{RS}}$（或等于）与 $q$ 闭式生成；实现须与正式标准表/公式一致。
\end{itemize}

\textbf{6.5.9.1 组跳}：
\begin{align}
c_{\mathrm{init}}&=\lfloor n_{\mathrm{ID}}^X/30\rfloor\\
f_{\mathrm{gh}}(n_S)&=\Bigl(\sum_{i=0}^{7}c(8n_S+i)\,2^i\Bigr)\bmod 30\\
f_{\mathrm{SS}}&=\lfloor n_{\mathrm{ID}}^X/16\rfloor\bmod 30\\
u&=(f_{\mathrm{gh}}+f_{\mathrm{SS}})\bmod 30
\end{align}

%----------------------------------------------------------------------
\section{本节核心详解}
%----------------------------------------------------------------------

\subsection{加扰}

发送：算 $c_{\mathrm{init}}$ $\rightarrow$ Gold $\rightarrow$ 异或。接收：同 $c_{\mathrm{init}}$ 再异或还原。每 TB 按起始帧重算，不跨 TB 保持加扰状态。

\subsection{SC-FDMA}

仅 6.3.3；校验整除后对每段独立 DFT。不做 IFFT/CP。

\subsection{ZC / 组跳}

同槽 $n_S$ 内 $u$ 不变；换槽先组跳再生成序列。$m{=}2$ 必须 ROM 查表。

%----------------------------------------------------------------------
\section{数据类型、长度与维度}
%----------------------------------------------------------------------

\begin{longtable}{L{3.5cm}L{2.5cm}L{8cm}}
\caption{关键接口维度\label{tab:6579-dim}} \\
\toprule
\textbf{变量} & \textbf{方向} & \textbf{说明} \\
\midrule
\texttt{bits\_in/out} & 入/出 & $N_{\mathrm{bit}}$ 比特加扰 \\
\texttt{c\_init} & 输入 & Gold 初始化 \\
\texttt{x[]/y[]} & 入/出 & $M_{\mathrm{symb}}$ 复数；DFT 长 $M$ \\
\texttt{u,v,α,m} & 输入 & ZC 参数 \\
\texttt{n\_S,n\_ID\_X} & 输入 & 组跳 \\
\bottomrule
\end{longtable}

%----------------------------------------------------------------------
\section{时序关系与触发条件}
%----------------------------------------------------------------------

\begin{longtable}{L{3.5cm}L{5.5cm}L{5.5cm}}
\caption{调用触发条件\label{tab:6579-trig}} \\
\toprule
\textbf{能力} & \textbf{进入条件} & \textbf{失败/回退} \\
\midrule
加扰 & 编码比特就绪，$c_{\mathrm{init}}$ 已定 & 空指针 $\rightarrow$ \texttt{ERR\_PARAM} \\
SC-FDMA & $M_{\mathrm{symb}}\bmod M{=}0$ & 不整除 $\rightarrow$ \texttt{ERR\_PARAM} \\
ZC $m{=}2$ & $u\in[0,29]$，表可用 & 误走闭式 $\rightarrow$ \texttt{ERR\_MODE} \\
组跳 & $n_S$、$n_{\mathrm{ID}}^X$ 有效 & 与加扰共用 Gold 偏移 $\rightarrow$ 不合格 \\
\bottomrule
\end{longtable}

%----------------------------------------------------------------------
\section{处理步骤}
%----------------------------------------------------------------------

\subsection{端到端链路 A：数据加扰 → 调制}

\textbf{Step 1}：按业务查表算 $c_{\mathrm{init}}$。\\
\textbf{Step 2}：\texttt{slb\_pn\_gold\_generate} 得 $c(i)$。\\
\textbf{Step 3}：$\tilde{t}_i=t_i\oplus c(i)$。\\
\textbf{Step 4}：交给 6.5.6。

\subsection{端到端链路 B：深覆盖 SC-FDMA}

\textbf{Step 1}：校验 $M_{\mathrm{symb}}\bmod M_{\mathrm{SC}}^{\mathrm{Total}}{=}0$。\\
\textbf{Step 2}：对每段 $l$ 做 $M$ 点 DFT 并乘 $1/\sqrt{M}$。\\
\textbf{Step 3}：输出 $y[]$ 给映射。

\subsection{端到端链路 C：DMRS}

\textbf{Step 1}：若使能组跳，算 $u$。\\
\textbf{Step 2}：$m{=}2$ 查表，否则按标准闭式生成 $r_{u,v}$。\\
\textbf{Step 3}：乘 $e^{j\alpha n}$；下游再乘 OCC/$β_{\mathrm{RS}}$。

%----------------------------------------------------------------------
\section{协议中允许本模块改变的参数}
%----------------------------------------------------------------------

\begin{longtable}{L{4cm}L{4cm}L{6cm}}
\caption{允许配置的参数\label{tab:6579-tunable}} \\
\toprule
\textbf{参数} & \textbf{来源} & \textbf{影响} \\
\midrule
加扰 $c_{\mathrm{init}}$ & 帧/$N_{\mathrm{G-PID}}$/业务 & 加扰图案 \\
$M_{\mathrm{SC}}^{\mathrm{Total}}$ & 带宽配置 & DFT 长度 \\
$u,v,\alpha,m$ & DMRS 调度 & ZC 序列 \\
$n_S,n_{\mathrm{ID}}^X$ & 时隙/业务 ID & 组跳 $u$ \\
\bottomrule
\end{longtable}

%----------------------------------------------------------------------
\section{链路调用对照}
%----------------------------------------------------------------------

\begin{longtable}{L{4cm}L{2.8cm}cc}
\caption{链路与 6.5.7/6.5.8/6.5.9 对照\label{tab:6579-link}} \\
\toprule
\textbf{链路} & \textbf{标准} & \textbf{6.5.7} & \textbf{6.5.8/9} \\
\midrule
第一/二类数据 & 6.2.5 & 加扰 & --- \\
系统广播 & 6.3.3.3 & $N_{\mathrm{ID}}^{\mathrm{TS}}$ & SC-FDMA \\
控制信息 & 6.3.3.4 & 510 & SC-FDMA \\
PHY 数据 & 6.3.3.5 & 高层 & SC-FDMA \\
DMRS & 6.3.3.2.2 & --- & ZC+组跳 \\
\bottomrule
\end{longtable}

%----------------------------------------------------------------------
\section{约束与实现要点}
%----------------------------------------------------------------------

\begin{enumerate}[nosep]
  \item \textbf{【先加扰后调制】}：6.5.6 数据输入必须是加扰比特；
  \item \textbf{【$c_{\mathrm{init}}$ 隔离】}：6.5.3 / 6.5.7 / 6.5.9.1 三类初始化不得混用或共享流偏移；
  \item \textbf{【DFT 整除】}：$M_{\mathrm{symb}}$ 必须被 $M_{\mathrm{SC}}^{\mathrm{Total}}$ 整除；
  \item \textbf{【$1/\sqrt{M}$】}：SC-FDMA 归一化不得删改；
  \item \textbf{【$m{=}2$ 查表】}：禁止对 $m{=}2$ 套用 ZC 闭式；
  \item \textbf{【组跳每槽更新】}：同槽 $u$ 固定，换槽先跳再生成；
  \item \textbf{【无跨 TB 加扰状态】}：每 TB 按起始帧重算 $c_{\mathrm{init}}$；
  \item \textbf{【无状态机】}：加扰/DFT/ZC 均按条件直接计算；不维护编码态/加扰态/调制态/SC-FDMA 态等 FSM 跳转表。
\end{enumerate}

%----------------------------------------------------------------------
\section{与标准其他章节的衔接}
%----------------------------------------------------------------------

\begin{itemize}[nosep]
  \item \textbf{6.5.2}：Gold 公共引擎；
  \item \textbf{6.5.6}：加扰后调制；调制后可选 SC-FDMA；
  \item \textbf{6.2.5}：常规数据仅 6.5.7，无 6.5.8；
  \item \textbf{6.3.3}：深覆盖完整链与 DMRS OCC/$β_{\mathrm{RS}}$。
\end{itemize}

\vfill
\begin{center}
\small\textcolor{gray}{精确参数与字段定义以 T/XS 10001-2025 正式标准为准；\\
本文档提供 6.5.7\textasciitilde 6.5.9 工程解读，供 SLB 实现与联调参考。}
\end{center}

\end{document}
